\subsection{Avalanche Ionization}

The energy shift is given by $\Delta E = \hbar\omega_0 \pm \hbar\omega_q$, where $\hbar\omega_q$ is the phonon energy and $\hbar\omega_0$ is the photon energy. Since the phonon energy is very small compared to the photon energy, we can assume $\Delta E \approx \hbar\omega_0$. Electrons can absorb several photons in sequence, increasing their velocity and their energy. Beyond a certain threshold, the energy of the free carriers exceeds the gap energy, and they can then collide with electrons of the valence band, causing impact ionization. This leaves two electrons in the conduction band, which increases the absorption efficiency and triggers further impact ionization events. This cascading effect is known as avalanche ionization. It combines free carrier absorption with impact ionization.

This process can repeat itself as long as the electric field is present and intense enough. The population of the conduction band is set by the competition between the photoionization rate and the avalanche rate. In dielectric materials such as quartz or fused silica, self-trapping and recombination occur on a time scale comparable to the duration of a femtosecond laser pulse. They must therefore be included in the population model, in the form of a system of two coupled populations~\cite{mao2004}:
\begin{align}
    \frac{dN}{dt} &= N W_{PI}(|\mathcal{E}|) + \frac{\sigma}{E_g} N |\mathcal{E}|^2 + N_{STE} W_{PI}(|\mathcal{E}_x|) - \frac{N}{\tau_x}, \\
    \frac{dN_{STE}}{dt} &= -N_{STE} W_{PI}(|\mathcal{E}_x|) + \frac{N}{\tau_x},
\end{align}
where $W_{PI}$ is the Keldysh ionization rate, $E_g$ is the gap energy, $\sigma = \frac{k_0 e^2}{n_0^2 \omega_0^2 \varepsilon_0 m_e} \frac{\omega_0 \tau_c}{1 + \omega_0^2 \tau_c^2}$ is the inverse Bremsstrahlung cross section derived from the Drude model, and $\tau_x$ is the characteristic self-trapping time.
